% !TeX root = ../main.tex

\section{定義・定理}

\subsection{基本的な定義}
\begin{comment}
	\begin{definition}
		% 定義を書く
	\end{definition}
	
	\begin{remark}
		% 定義に関する注意を書く
	\end{remark}
	
	\subsection{基本的な定理}
	
	\begin{theorem}
		% 定理を書く
	\end{theorem}
	
	\begin{lemma}
		% 必要なら補題を書く
	\end{lemma}
	
	\begin{corollary}
		% 必要なら系を書く
	\end{corollary}
\end{comment}

\begin{theorem}[収束列の部分列]
	数列 $\{a_n\}_{n\in\mathbb{N}}$ が $\alpha\in\mathbb{R}$ に収束するとする。
	また，$\{a_{n_k}\}_{k\in\mathbb{N}}$ を $\{a_n\}$ の部分列とする。
	このとき，
	\[
		\lim_{k\to\infty}a_{n_k}=\alpha
	\]
	である。
\end{theorem}
\begin{proof}
	数列 $\{a_n\}$ が $\alpha$ に収束すると仮定する。
	すなわち，任意の $\epsilon>0$ に対して，ある自然数 $N$ が存在して，
	$n\geq N$ ならば
	\[
		|a_n-\alpha|<\epsilon
	\]
	が成り立つ。

	いま，$\{a_{n_k}\}_{k\in\mathbb{N}}$ を $\{a_n\}$ の部分列とする。
	部分列であるから，添字の列
	\[
		n_1<n_2<n_3<\cdots
	\]
	は自然数からなる狭義単調増加列である。

	このとき，すべての $k\in\mathbb{N}$ に対して
	\[
		n_k\geq k
	\]
	が成り立つ。

	実際，$n_1\geq 1$ であり，添字は少なくとも $1$ ずつ増えるので，
	\[
		n_k\geq k
	\]
	である。

	ここで，任意に $\epsilon>0$ を取る。
	$a_n\to\alpha$ より，ある自然数 $N$ が存在して，
	$n\geq N$ ならば
	\[
		|a_n-\alpha|<\epsilon
	\]
	が成り立つ。

	そこで，部分列の番号について
	\[
		K=N
	\]
	とおく。

	このとき，$k\geq K$ ならば
	\[
		k\geq N
	\]
	である。さらに $n_k\geq k$ だから，
	\[
		n_k\geq k\geq N
	\]
	となる。

	したがって，元の数列の収束性より，
	\[
		|a_{n_k}-\alpha|<\epsilon
	\]
	が成り立つ。

	つまり，任意の $\epsilon>0$ に対して，ある自然数 $K$ が存在して，
	$k\geq K$ ならば
	\[
		|a_{n_k}-\alpha|<\epsilon
	\]
	である。

	よって，数列 $\{a_{n_k}\}$ は $\alpha$ に収束する。
	すなわち，
	\[
		\lim_{k\to\infty}a_{n_k}=\alpha
	\]
	である。
\end{proof}